% БҮЛЭГ 1 — АСУУДЛЫН ТОМЬЁОЛОЛ (УДИРТГАЛ)

\chapter{Удиртгал}
\label{Chapter1}

Энэхүү бүлэгт дипломын судалгааны үндэслэл, шийдэж буй асуудлын тодорхой
томьёолол, SMART зарчмын дагуу томьёолсон зорилго ба зорилтууд, ажлын
таамаглал, хамрах хүрээ болон шинжлэх ухаан, нийгэмд оруулах хувь нэмрийг
нарийвчлан танилцуулна.

%-------------------------------------------------------------------------------
\section{Судалгааны үндэслэл}

Интернэтийн хэрэглэгчдийн тоо дэлхий даяар урьд өмнө байгаагүй хурдацтай
нэмэгдэж байгаа бөгөөд цахим мэдээллийн орчин нийгмийн харилцааны хэлбэрийг
бүрэн өөрчилсөн. DataReportal (2025)-ийн мэдээллийн дагуу Монгол Улсын
интернэт хэрэглэгчдийн тоо нийт хүн амын 83 гаруй хувийг эзэлж байгаа бөгөөд
Facebook, X (Twitter), TikTok зэрэг платформуудын хэрэглэгчид жил бүр хурдацтай
өсч байна. Цахим орчинд тархаж буй мэдээлэл, харилцааны хэмжээ нэмэгдэхийн
хэрээр хэрэглэгчид, ялангуяа өсвөр насныхан, эрүүл бус, хортой агуулгад
өртөх эрсдэл ч нэгэн зэрэг нэмэгдэж байна.

Монголын цахим орчинд --- тухайлбал news.mn зэрэг мэдээний сайтын сэтгэгдлийн
хэсэг, Facebook-ийн нийтийн хуудас болон бүлгүүд, X (Twitter)-ийн нийтийн
нийтлэлүүдэд --- кибер дарамт, үзэн ядалтын хэл, хорлолтой дезинформаци, хэт
аггрессив хандлага зэрэг олон хэлбэрийн сөрөг агуулга тархаж байна. Олон
улсын судалгаанууд \cite{kowalski2014,vogels2022} кибер дарамт нь өсвөр насны
хүүхдийн сэтгэл зүйн тогтвортой байдал, нийгмийн харилцаа, академик
гүйцэтгэлд ноцтой сөргөөр нөлөөлдгийг нотолсон. Монгол Улсад энэ асуудалд
хариулах техникийн шийдэл, тухайлбал автомат контент шүүлтийн систем,
хөгжлийн түвшин хангалтгүй байгааг тэмдэглэх нь зүйтэй.

Гэхдээ аливаа сөрөг агуулгыг цогцоор нь хориглох нь нийгмийн хувьд буруу
шийдэл болно. Бүтээлч шүүмжлэл, асуудалд суурилсан шинжилгээ, бодлогыг
эсэргүүцэх тайлбар зэрэг нь ардчилсан нийгэм, иргэний соёлын чухал суурь
болдог. ``Positive-only'' орчин нь нийгмийн бодит байдлаас тасарсан хиймэл
дискурсийг бий болгох аюултай бөгөөд үзэл бодлоо чөлөөтэй илэрхийлэх нь
үндсэн хүний эрхэд хамаарна. Тиймд автомат шүүлтийн систем нь зөвхөн бүх
сөрөг агуулгыг устгах зорилготой байж болохгүй, харин хоёр эсрэг туйлын
агуулгыг --- нэг нь хориглох, нөгөөг нь хамгаалах --- нарийн ялгаж чаддаг
байх шаардлагатай.

%-------------------------------------------------------------------------------
\section{Асуудлын томьёолол}

Монгол хэлний тусгайлсан NLP судалгааны хүрээнд Toxic болон Constructive
агуулгыг ялгах систем одоогоор боловсруулагдаагүй байгаа нь судалгааны
цоорхойг бүрдүүлж байна. Өмнөх ажлуудын дийлэнх нь гурван ангилалт
(Positive / Negative / Neutral) сентимент шинжилгээгээр хязгаарлагдаж,
``Negative'' ангилалын дотоодод байх хоёр эсрэг шинж чанарыг --- устгах
шаардлагатай хортой агуулга ба хамгаалах шаардлагатай бүтээлч шүүмжлэл ---
ялган тодорхойлоогүй байна. Энэхүү судалгааны ажил тэрхүү цоорхойг нөхөх
зорилготой бөгөөд нийгмийн болон техникийн ач холбогдлоороо тодорхой хувь
нэмэр оруулна.

%-------------------------------------------------------------------------------
\section{Судалгааны зорилго}

Монгол хэл дээрх цахим орчны агуулгыг олон ангилалтаар шүүх, тухайлбал Toxic
Negative болон Constructive Criticism-ийг нарийн ялгаж чаддаг хиймэл оюунд суурилсан
NLP системийг боловсруулах нь энэхүү судалгааны гол зорилго юм. Тус систем нь
Монгол хэлний agglutinative морфологийн онцлогт тохирсон өгөгдлийн
боловсруулалтын дамжлага (pipeline), MN-BERT загварын нарийн тохируулга
(fine-tuning), болон Gradio-д суурилсан хэрэглэгчийн интерфейсийг нэгтгэнэ.

Зорилгыг SMART зарчмын дагуу томьёолбол:

\begin{itemize}
  \item \textbf{Specific (тодорхой):} 4 ангилалтай (Positive, Neutral, Toxic
    Negative, Constructive Criticism) Монгол кирилл текстийн ангилагч
    системийг хоёр шатлалт архитектуртай зохион бүтээх.
  \item \textbf{Measurable (хэмжигдэхүйц):} End-to-end macro F1 $\geq 0.80$,
    нэгдсэн accuracy $\geq 90\%$ (guideline Ш-1) зорилтоор хэмжих.
  \item \textbf{Achievable (биелэгдэхүйц):} MN-BERT-ийн pre-trained жинг
    суурь болгон fine-tune хийж, SMOTE болон жингэсэн loss-оор ангиллын
    тэнцвэргүй байдлыг шийдвэрлэх.
  \item \textbf{Relevant (хамааралтай):} Монгол хэлний NLP-ийн одоогийн
    ``Positive/Neutral/Negative'' гурван ангилалт суурь дээр Toxic vs
    Constructive ялгалт нэмэлтээр оруулах цоорхойтой шууд холбоотой.
  \item \textbf{Time-bound (хугацаатай):} 2025 оны 9 сараас 2026 оны 5 сар
    хүртэлх дипломын ажлын хугацаанд хэрэгжүүлэх.
\end{itemize}

%-------------------------------------------------------------------------------
\section{Судалгааны зорилтууд}

Дипломын ажлын хүрээнд дараах зорилтуудыг дараалсан байдлаар хэрэгжүүлнэ.

\begin{enumerate}
  \item news.mn болон Facebook нийтийн хуудас, бүлгүүдээс Монгол хэл дээрх
        сэтгэгдэл болон нийтлэлүүдийг цуглуулах, цэвэрлэх, стандартжуулах
        ажлыг гүйцэтгэнэ.

  \item Цуглуулсан өгөгдлийг Positive, Neutral, Toxic Negative, Constructive
        Criticism гэсэн дөрвөн ангилалд гар аргаар шошголох (manual
        annotation) бөгөөд аннотаторуудын хоорондын нийцлийг Cohen's Kappa
        коэффициентээр хэмжинэ.

  \item Монгол хэлний agglutinative онцлогт тохирсон урьдчилсан боловсруулалтын
        (preprocessing) pipeline --- токенчлол, стоп үг хасах, нормчлол ---
        боловсруулна.

  \item Хоёр шатлалт ангиллын архитектур зохион байгуулна: Stage~1 нь
        Positive/Neutral/Negative ангилал; Stage~2 нь Negative агуулгыг Toxic
        ба Constructive гэж нарийвчилсан ялгалт.

  \item MN-BERT загварыг fine-tuning хийж Logistic Regression, Naive Bayes,
        SVM зэрэг суурь загваруудтай харьцуулан статистик аргаар үнэлнэ.

  \item Accuracy, Precision, Recall, $F_1$-score, Confusion Matrix зэрэг
        үнэлгээний үзүүлэлтүүдэд тулгуурлан загваруудын гүйцэтгэлийг
        тодорхойлно.

  \item Gradio framework ашиглан хэрэглэгчийн интерфейстэй прототип
        системийг хэрэгжүүлж, туршилтын орчинд үнэлнэ.
\end{enumerate}

%-------------------------------------------------------------------------------
\section{Таамаглал}

Судалгааны гол таамаглал нь: \textit{``Монгол хэлний agglutinative онцлогт
тохирсон preprocessing pipeline дээр суурилан, MN-BERT pre-trained загварыг
хоёр шатлалт (two-stage) архитектурт fine-tune хийж ашиглавал, TF-IDF
суурьтай уламжлалт машин сургалтын baseline (LR, NB, SVM) загваруудаас
дээгүүр macro F1 гүйцэтгэл үзүүлж, Toxic Negative болон Constructive
Criticism-ийг найдвартай ялгах боломжтой.''}

Дэд таамаглалууд:
\begin{itemize}
  \item H$_1$: MN-BERT fine-tuned загвар нь TF-IDF + LR/NB/SVM-ээс macro F1
    хэмжүүрээр дор хаяж 5 хувийн цэг дээгүүр гүйцэтгэл үзүүлнэ.
  \item H$_2$: Ангиллын тэнцвэржүүлэлт (SMOTE baseline-д, жингэсэн loss
    MN-BERT-д) нь цөөнх ангиллын F1-г 10 хувиар сайжруулна.
  \item H$_3$: Хоёр шатлалт архитектур нь Toxic vs Constructive-ийг flat
    4-ангиллын загвараас илүү нарийвчилсан тооцоолох pipeline өгнө (хэдий
    каскад алдааны эрсдэлтэй ч).
\end{itemize}

%-------------------------------------------------------------------------------
\section{Судалгааны хамрах хүрээ ба хязгаарлалт}

Энэхүү судалгаа нь дараах хэмжээст орон зайд явагдана. Хэлний хувьд зөвхөн
Монгол кирилл бичгийн текст хамрагдана. Латин үсгийн бичвэр, монгол-англи
холимог (code-switched) текст, романчилсан Монгол бичлэг (``bn'', ``yu ve'',
``zgr'' зэрэг), болон бусад кирилл бус бичиглэлийн хэлбэрүүд датасетаас
тусгайлан хасагдана. Энэхүү шийдэл нь MN-BERT-ийн кирилл корпус дээр суурилсан
токенчлолын тогтвортой байдлыг хангаж, аннотацийн стандартыг нэгтгэхэд
шаардлагатай.

Платформын хувьд news.mn сэтгэгдлийн хэсэг болон Facebook нийтийн хуудас,
нийтийн нийтлэлүүд хамрагдана. Хугацааны хувьд 2024--2025 оны нийтлэлүүдийг
ашиглана. Контентийн хувьд зөвхөн текстэн өгөгдлийг анализад оруулах бөгөөд
зураг, видео, эможи зэргийг анхдагч загварт оруулахгүй --- эдгээрийг ирээдүйн
судалгааны чиглэл болгон тодорхойлно. Системийн хувьд прототип загвар
боловсруулахаар тооцох бөгөөд production орчинд байршуулах нь энэхүү ажлын
хүрээнд багтахгүй.

%-------------------------------------------------------------------------------
\section{Ажлын шинэлэг тал ба ач холбогдол}

Дипломын ажлын шинэлэг тал нь Монгол хэлний цахим орчинд чиглэсэн \textit{анхны}
хоёр шатлалт Toxic/Constructive ангиллын системийг нэвтрүүлж буйд оршино.
Үүнээс гадна: (1)~Монгол хэлний agglutinative онцлогт тохирсон нэгдсэн
preprocessing pipeline, (2)~MN-BERT-ийн fine-tuning-ыг FocalLoss болон
cosine scheduler-тэй хослуулсан туршилт, (3)~inter-annotator agreement-ыг
ил тод хэмжиж баримтжуулсан өгөгдлийн багц --- гурав нь салбарын одоогийн
төлөв байдалд тодорхой хувь нэмэр болно.

Нийгмийн ач холбогдлын хувьд уг систем нь Монгол цахим орон зайд сэтгэл
хор, дарамт, үзэн ядалт зэрэг сөрөг агуулгыг хянан илрүүлэх хэрэгсэл болох
бөгөөд хамгаалагдсан хууль ёсны шүүмжлэл хэлбэрийн үзэл бодлыг автоматаар
таслан зогсоохоос урьдчилан сэргийлэх ёс зүйн хэв загвар бүрдүүлнэ.

%-------------------------------------------------------------------------------
\section{Бүлгийн дүгнэлт}

Монгол хэлний цахим орчинд хортой болон бүтээлч агуулгыг автоматаар ялган
шүүх асуудал нь зөвхөн техникийн биш, нийгэм, иргэний болон хүний эрхийн
чухал ач холбогдолтой. Өмнөх судалгааны ажлуудад энэ ялгалт хийгдэж
байгаагүй нь тодорхой судалгааны цоорхойг бий болгосон. Энэхүү бүлэгт
асуудлын томьёолол, SMART зорилго, гурван дэд таамаглал, хамрах хүрээ,
шинэлэг тал болон нийгмийн ач холбогдлыг тодорхой тогтоов. Дараагийн
бүлэгт (бүлэг~\ref{Chapter2}) NLP, мэдрэмжийн шинжилгээ, BERT архитектур
болон холбогдох ажлуудын онолын суурийг нарийвчлан авч үзнэ.
